% 注：1. pdfcover的内容根据学校要求自定义（参考cover_file文件夹）。2. 这里代码写得不是很完备，第二个参数只影响封面，其他内容如页眉不受影响。还需在下面代码进一步调整。这里还需完善。为简单期间，下面三个选项供选择：
\documentclass[unicode,master]{scutthesis} % 硕士草稿封面，使用正式封面时注释该行
% \documentclass[unicode]{scutthesis} %  博士草稿封面，使用正式封面时注释该行
% \documentclass[unicode,pdfcover]{scutthesis} %   论文正式封面（不使用选项来区分硕博士，而是大家制作自己的封面文件thesis_cover.pdf，盲审时删除必要的封面内容即可），使用草稿封面时注释该行

\usepackage{array,longtable,graphicx} 
\usepackage{anyfontsize} %消除字体警告
\usepackage{enumitem}
\usepackage{bm}% 粗体数学符号
%%%%%%%%%%%%%%%%%%%%%%%%%%%%%%%%%%%%%%%%%%%%%%%%%%%%%%%%%%%%%%%%%%%%%%%%%%%%%%%%%%
%编译范围

\includeonly{section/1_绪论}
% \includeonly{section/2_预备知识}
% \includeonly{section/3_正文}
% \includeonly{section/4_正文}
% \includeonly{section/5_正文}


%参考文献设置
\usepackage[backend=biber,style=gb7714-2015,gbalign=gb7714-2015,gbpub=false,gbnamefmt=lowercase,gbpunctin=false]{biblatex}  % 添加gbpuctin=false用于去除//
% 新添一行代码用于去掉会议前的In
\DefineBibliographyStrings{english}{in={}}


\addbibresource[location=local]{reference.bib} 
%页眉页脚设置
\usepackage{fancyhdr}
\usepackage{listings}
\usepackage{xunicode}
\renewcommand{\lstlistingname}{列表}
\pagestyle{fancy}
\fancyfoot[C]{\headfont\thepage}
\renewcommand{\chaptermark}[1]{\markboth{\chaptername\ #1}{}}
\renewcommand{\sectionmark}[1]{\markright{\thesection\ #1}}
\fancyhead[RE]{}
\fancyhead[RO]{}
\fancyhead[LE]{}
\fancyhead[LO]{}
\fancyhead[CO]{\headfont{\leftmark}}
\fancyhead[CE]{\headfont{华南理工大学硕士学位论文}}% 博士自行修改
\renewcommand{\headrulewidth}{1.5pt}
\renewcommand{\footrulewidth}{0pt}
%%%%%%%%%%%%%%%%%%%%%%%%%%%%%%%%%%%%%%%%%%%%%%%%%%%%%%%%%%%%%%%%%%%%%%%%%%%%%%%%%%
\usepackage[unicode=true,bookmarks=true,bookmarksnumbered=true,bookmarksopen=false,breaklinks=false,pdfborder={0 0 1},backref=false,colorlinks=true]{hyperref}
%  \hypersetup{...}主要影响生成的 PDF 文件的元数据（PDF属性）和超链接样式，不会影响论文正文内容的显示。
\hypersetup{pdftitle={基于观测输入的机器人模仿学习研究},
	pdfauthor={姓名}, % 可能会泄露作者信息，盲审时请注意
	pdfsubject={华南理工大学硕士学位论文}, % 博士自行修改
	pdfkeywords={机器人;模仿学习;观测输入;深度学习},
%%		linkcolor=black, anchorcolor=black, citecolor=black, filecolor=black, menucolor=black, urlcolor=black, pdfstartview=FitH}% 黑白，提交版
	linkcolor=blue, anchorcolor=black, citecolor=red, filecolor=magenta, menucolor=red, urlcolor=magenta, pdfstartview=FitH}% 彩色

\makeatletter
%%%%%%%%%%%%%%%%%%%%%%%%%%%%%% LyX specific LaTeX commands.
\providecommand{\LyX}{\texorpdfstring%
	{L\kern-.1667em\lower.25em\hbox{Y}\kern-.125emX\@}
	{LyX}}
%% Because html converters don't know tabularnewline
\providecommand{\tabularnewline}{\\}
\makeatother
\begin{document}
	%%%%%%%%%%%%%草稿封面设置%%%%%%%%%%%%%使用“正式封面”时不需要理会这部分
	\title{基于观测输入的机器人模仿学习研究}
	\author{许文晋}
	\supervisor{指导教师：杨辰光\ 教授}
	\institute{华南理工大学}
	\date{2026年4月}
	%%%%%%%%%%%%%%%%%%%%%%%%%%%%%%%%%%%%%
	\maketitle	
	\frontmatter	%此后为罗马数字页码，页面类型为plain
	\chapter{摘\texorpdfstring{\quad}{}要}

随着全球产业变革的推进，机器人的应用正从结构化的工业环境转向高度非结构化的居家场景。然而，现有的机器人模仿学习系统仍面对专家示教数据匮乏、非结构化环境约束复杂以及现实多模态任务决策困难等挑战。因此，研究基于观测输入的机器人模仿学习方法，对于推动机器人智能化水平的提升具有重要意义。

首先，针对模仿学习系统中专家数据稀缺的难题，本文基于机器人书法任务提出一种基于图像观测输入的书法技能模仿学习框架，实现从海量汉字书法图像中学习书法技能。该框架利用八邻域优先级算法实现笔画拓扑骨架的精准分割，并结合模糊高斯混合模型解决高曲率笔画的非线性拟合难题。通过引入偏差分析算法，从像素中重构表征笔画粗细变化的协方差序列，并结合动态运动原语实现了书法技能在不同空间尺度下的自适应泛化。搭建了基于艾利特机器人的汉字书写实验平台，在六个书写评价指标的平均提升了10.16\%，并测试了多个不同汉字的书写。

其次，针对模仿学习系统在非结构化环境中难以适应异构约束的问题，本文提出了基于异构约束观测输入的模仿学习网络。该网络通过编码空间障碍物位置或运动学限制等约束观测信息，利用零卷积初始化与残差注入机制，在冻结的扩散策略骨干中引入引导偏置，将物理约束信息注入到策略网络。仿真实验论证了不同物理约束在特征空间中的解耦特性，实现了多约束条件的零样本组合与即插即用式适配。在具有障碍物的仿真任务中，该方法达到了 90.5\% 的有效成功率，大幅增强了机器人复杂操作的安全性。

最后，为了提升模仿学习系统在现实环境中处理复合约束，动态约束和多任务泛化的能力，本文构建了基于多模态观测输入的模仿学习系统。利用自监督学习策略构建语音去噪网络，实现了基于语音指令的约束模块动态调度，在真实避障场景、运动学约束场景和复合约束场景中的成功率均在80\%以上。集成了实时目标检测模块提取障碍物几何特征，构建了闭环感知反馈回路，在长程动态避障分拣场景中的有效成功率达到 100\%。设计了基于对比学习的潜动作编码器，构建了多构型统一的动作表征，在后训练实现一倍的收敛速度提升。通过引入动作后处理机制，解决了推理频率与执行频率不匹配导致的运动抖动，使叠毛巾任务的成功率提升至 72\%，抓取任务成功率达到 90\%，提升了机器人在执行复杂任务时的运动稳定性与多任务泛化能力。

\keywordsCN{机器人模仿学习；高斯混合模型；动态运动原语；扩散策略；视觉-语言-动作模型}

\chapter{Abstract}
With the advancement of global industrial transformation, robots are progressively transitioning from structured industrial environments to highly unstructured domestic service scenarios. Nevertheless, existing robot imitation learning systems continue to face significant challenges, including the scarcity of expert demonstration data, the complexity of unstructured environmental constraints, and the difficulty of multi-modal task decision-making in real-world scenarios. Therefore, research on robot imitation learning methods based on observational inputs is of considerable significance for advancing robot intelligence and facilitating the large-scale deployment of embodied intelligence.

Firstly, to address the challenge of scarce expert data in imitation learning systems, this dissertation proposes an image observation-based calligraphy skill imitation learning framework leveraging robot calligraphy tasks, enabling skill acquisition from large-scale Chinese character calligraphy images. The framework employs an eight-neighbor priority algorithm to achieve precise segmentation of stroke topological skeletons, and incorporates a fuzzy Gaussian mixture model to address the nonlinear fitting challenges inherent in high-curvature strokes. Furthermore, a deviation analysis algorithm is introduced to reconstruct the covariance sequence characterizing stroke width variations from pixel data. This covariance information is then integrated with dynamic movement primitives to achieve adaptive generalization of calligraphy skills across different spatial scales. An experimental platform based on an Elite robot is established for Chinese character writing tasks. The proposed framework demonstrates an average improvement of 10.16\% across six writing evaluation metrics, and successful generalization is validated through testing on multiple different Chinese characters.

Secondly, to address the problem of imitation learning systems struggling to adapt to heterogeneous constraints in unstructured environments, this dissertation proposes an imitation learning network based on heterogeneous constraint observation inputs. The network encodes constraint observation information such as spatial obstacle positions or kinematic limitations, and utilizes zero- convolution initialization with residual injection mechanisms to introduce guiding biases into a frozen diffusion policy backbone, thereby injecting physical constraint information into the policy network. Simulation experiments demonstrate the decoupling characteristics of different physical constraints in the feature space, achieving zero-shot composition and plug-and-play adaptation of multiple constraint conditions. In obstacle-avoidance simulation tasks, this method achieves an effective success rate of 90.5\%, significantly enhancing robot operation safety in complex scenarios.

Finally, to enhance the capability of imitation learning systems to handle composite constraints, dynamic constraints, and multi-task generalization in real-world environments, this dissertation constructs an imitation learning system based on multi-modal observation inputs. A speech denoising network is built using a self-supervised learning strategy to achieve dynamic scheduling of constraint modules based on voice commands, with success rates exceeding 80\% in real obstacle-avoidance scenarios, kinematic constraint scenarios, and composite constraint scenarios. A real-time object detection module is integrated to extract obstacle geometric features and construct a closed-loop perceptual feedback loop, achieving an effective success rate of 100\% in long-horizon dynamic obstacle-avoidance sorting tasks. A latent action encoder based on contrastive learning is designed to construct a unified action representation across multiple configurations, achieving a two-fold improvement in convergence speed during post-training. By introducing an action post-processing mechanism, the motion jitter caused by the mismatch between inference frequency and execution frequency is resolved, raising the success rate of towel-folding tasks to 72\% and achieving a 90\% success rate in grasping tasks, thereby enhancing motion stability and multi-task generalization capabilities of robots when executing complex tasks.

\keywordsEN{Robotic Imitation Learning; Gaussian Mixture Model; Dynamic Movement Primitives; Diffusion Policy; Vision-Language-Action Models} % 中英文摘要
	%%%%%%%%%%%%%%%%%%%%%%%%%%%%%%%%%%%%%%%%%%%%%%%%
	% 目录、表格目录、插图目录这几个字本身的大纲级别是一级的，即和章名有相同的字号字体。目录表的内容通过titletoc宏包在。cls文件设置了。
	%\cleardoublepage % pdfbookmark可能需要这一条才能正常工作
	\pdfbookmark{\contentsname}{toc} %为目录添加pdf文件书签
	\tableofcontents	%目录

	
	\mainmatter %此后为阿拉伯数字页码
	
    %%%%%%%%%%%%%%%%%%%%%%%%%%%%%%%%%%%%%%%%%%%%%%页眉页脚设置  
    \fancypagestyle{plain}{
    	\pagestyle{fancy}
    }	% 每章的第一页会默认使用plain，没有页眉。通过重定义plain为fancy解决
    \pagestyle{fancy}	%设置页眉页脚为fancy
    %%%%%%%%%%%%%%%%%%%%%%%%%%%%%%%%%%%%%%%%%%%%%%分章节，结合导言区的\includeonly命令可仅编译部分章节，编译时不用切换界面，直接在相应章节编译即可。
	\chapter{绪论}

\section{研究背景与意义}

\section{国内外研究现状}

\section{论文研究内容}

\section{论文结构安排}
	\chapter{预备知识}
\label{chap:preliminary}

\section{高斯混合模型与动态运动原语基础}

\subsection{高斯混合模型 (GMM) 与高斯混合回归 (GMR)}

\subsection{动态运动原语 (DMP) 的二阶动态系统理论}

\section{生成式扩散策略 (Diffusion Policy)}

\subsection{扩散概率模型的正向加噪与反向去噪过程}

\subsection{视觉-本体条件引导机制 (Conditioning)}

\section{语音信号处理与注意力机制基础}

\subsection{时频域转换与复数谱特征}

\subsection{多头注意力机制与特征提取}

\section{本章小结}
	\chapter{基于图像观测输入的书法技能模仿学习框架}

\section{引言}

\section{基于八邻域优先级矩阵的笔画分割}

\section{基于模糊高斯混合模型的动态编码}

\section{基于动态运动原语的自适应空间泛化}

\section{实验对比与分析}

\section{本章小结}
	\chapter{基于异构约束观测输入的模仿学习网络}

\section{引言}

\section{基于扩散模型的模仿学习策略}

\section{面向异构约束的残差注入架构设计}

\section{零样本多约束组合机制}

\section{实验对比与分析}

\section{本章小结}
	\chapter{基于多模态观测输入的机器人系统}

\section{引言}

\section{复杂环境自监督语音降噪}

\section{基于视听融合的多模态约束机器人系统}

\section{基于动态视觉约束感知的闭环控制系统}

\section{基于多模态大模型的通用机器人策略}

\section{本章小结}

	\backmatter %章节不编号但页码继续
	%%%%%%%%%%%%%%%%%%%%%%%%%%%%%%%%%%%%%%%%%%%%%%%%%%%%%%%%%%%%%%    微调，使得后续章节的页眉不带章号
	\renewcommand{\chaptermark}[1]{\markboth{#1}{}}
	%%%%%%%%%%%%%%%%%%%%%%%%%%%%%%%%%%%%%%%%%%%%%%%%%%%%%%%%%%%%%%
	\chapter{总结与展望}
\section*{一、论文工作总结}
随着具身智能技术的飞速发展，机器人正加速从结构化的工业生产线走向非结构化、任务高度并行的居家与社会服务场景。然而，真实物理环境的不确定性、专家示教数据的匮乏以及任务指令的语义模糊性，对机器人模仿学习系统的感知深度与执行自主性提出了严峻挑战。本文针对这些痛点，深入开展了基于观测输入的机器人模仿学习研究，构建了从静态图像技能逆向获取到异构约束实时感知，再到通用多模态决策的完整技术体系。本文的主要工作总结如下：

(1) 针对书法等精密技能示教中第一人称视频稀缺且静态图像丢失动态信息的难题，本文设计并实现了一种基于图像观测输入的书法技能模仿学习框架。该框架通过设计八邻域优先级算法实现了对汉字笔画拓扑骨架的精准分割，并引入结合主动曲线轴机制的模糊高斯混合模型 (FGMM)，利用非线性轴线替代传统的线性主轴，解决了高曲率笔画在时序编码过程中的拟合难题。通过偏差分析 (DA) 算法，系统成功从像素分布中重构出表征笔画粗细变化的协方差序列，并利用高度映射机制将其转化为毛笔的法向控制指令。结合动态运动原语 (DMP) 对轨迹进行建模，实现了书法技能在不同空间尺度下的自适应泛化，确保了机器人在不同书写介质上的几何结构完整性与物理复现保真度。

(2) 为了解决生成式扩散策略在非结构化环境下难以实时遵循复杂环境约束的问题，本文设计了面向异构约束观测输入的模块化残差学习网络 ConstrainNet。该网络通过编码空间避障或运动学限制等异构观测信息，利用零卷积初始化与残差注入机制，在冻结的扩散策略骨干网络中引入引导偏置。这种设计确保了系统在训练初期能够保持预训练的操作先验，避免了全参数微调带来的灾难性遗忘。实验论证了空间约束与运动学约束在特征空间中的解耦特性，实现了多约束条件的零样本组合与即插即用式适配。线性探测与可视化分析表明，该网络成功在特征空间中重构了避障所需的几何语义，使机器人能够从被动的反应式滤波转向主动的全局轨迹规划，显著提升了作业安全性。

(3) 为了提升机器人系统在动态交互环境中的自主感知与决策能力，本文构建了基于多模态观测输入的闭环模仿学习系统。首先，利用自监督学习策略构建了语音去噪网络，通过最小化两组带噪信号间的差异，在无需干净标签的前提下实现了对环境指令的特征重建，并将其作为高层意图开关触发 ConstrainNet 模块调度。其次，通过集成 YOLOv8s 目标检测模型，系统能够直接从原始图像流中实时提取障碍物的四维边界框信息，并将其转化为策略网络的连续约束输入，实现了从被动指令触发向主动闭环感知的演进。在长时序分拣任务实验中，该系统展示了规避动态障碍物的能力，验证了残差注入机制在处理非稳态工况时的实时性与可靠性。

(4) 针对大规模机器人预训练中动作空间异构性的挑战，本文开发了端到端训练的通用视觉-语言-动作 (VLA) 模仿学习系统。系统设计了基于对比学习的潜动作编码器，利用监督对比损失与均匀性损失将不同构型机器人的观测-动作对映射至统一的潜空间，为跨任务、跨平台的技能迁移提供了共享表征。在大规模多模态数据集上对 VLA 基模进行预训练与后训练微调，验证了其在处理复杂语义指令时的泛化性能。针对真机部署中推理频率与执行频率不匹配导致的运动抖动问题，引入了基于三次多项式插值与滑动窗口滤波的后处理平滑机制。实验结果表明，该机制有效抑制了时序预测中的高频噪声，确保了机器人在执行叠衣、抓取等精细化家务任务时的运动连续性与稳定性。

\section*{二、未来展望}
针对本文已有研究的不足，对未来工作提出以下几点展望:

(1) 书法技能模仿学习框架存在以下局限性。框架假设笔画宽度一阶决定因素为毛笔法向按压力度，将书写速度引起的墨水渗透效应视为二阶残差，速度与压力的完全解耦在仅依赖静态图像输入的条件下不可能实现，未来可引入动态书写视频或加装触觉传感器来突破此约束；笔画分割算法基于"书写方向在交叉处不突变"假设，在行书草书多笔画缠绕等极端场景下失效，且仅在艾利特EC66单一平台上测试了有限楷体汉字。未来可在适应更多风格汉字的方向发展该框架。

(2) ConstrainNet框架的设计定位是软约束引导机制，提供特征层面的动作分布偏移而非可证明的硬安全保证。约束观测输入为低维代理表征（2D边界框或关节空间表达），深度方向避障隐式依赖策略网络的视觉泛化能力；对需绝对安全保障的场景，建议结合控制屏障函数等方法。在解决组合约束冲突时当前采取重采样方案，未来可从训练层面在损失函数中加入约束表征空间的正则化项。

(3) VLA系统的后处理机制在高速动态场景下存在适用边界。固定窗口大小（$W=5$）的滑动平均滤波器引入约200ms群延迟，在准静态任务（如叠毛巾）中可接受，但在末端速度超过1.0m/s的动态场景下将引入20cm以上物理偏差。未来可采用自适应窗口策略或BID、RTC等异步推理架构，使动作预测与执行解耦。

(4) VLA跨构型迁移实验目前仅在星海图R1 Lite单一新构型上验证，更大规模跨平台迁移能力仍需检验。潜动作编码器的对比学习损失未显式建模动力学参数，虽是为保持潜空间构型无关性的主动设计，但动力学差异可能对精确力控任务产生潜在影响。未来可引入如力信息以及世界模型等方法填补动力学建模的缺失。 %结论
	 %%%%%%%%%%%%%%%%%%%%%%%%%%%%%%%%%%%%%%%%%%%%%% bibtex参考文献设置  （原版）
%%	\bibliographystyle{scutthesis}
%%	\bibliography{F:/MyLibrary}
	%%%%%%%%%%%%%%%%%%%%%%%%%%%%%%%%%%%%%%%%%%%%%%
	%%%%%%%%%%%%%%%%%%%%%%%%%%%%%%%%%%%%%%%%%%%%%% biber参考文献设置	
	%\renewcommand*{\bibfont}{\refbodyfont}			% 设置文献著录字号比正文小一号（五号），需要小四号请注释该行. % 不推荐使用small，而是使用cls文件中精确定义了的字号。
	\phantomsection % “目录”中的链接能正确跳转，需要添加 \phantomsection 否则点击参考文献会跳转到结论
	\addcontentsline{toc}{chapter}{参考文献}	%目录中添加参考文献
	\printbibliography	% 参考文献著录
	\include{section/7_研究成果} %成果
	\chapter{致\texorpdfstring{\quad}{}谢}

回首在华园的七年时光，心中满是感激。

首先谨向我的导师杨辰光教授致以最诚挚的谢意。杨老师性格儒雅随和，治学严谨，在我心中是一位高尚的学者。每当我在研究中遇到困惑，杨老师总是耐心地给予指导，其渊博的学识让我受益匪浅。更令我感念的是，杨老师待人接物始终宽厚温和，让我学会如何以善意和包容对待身边的人。感谢实验室的同门李沐、杰彬、宇博、正聪、深元、思睿、雅静和志毅，和你们一起谈天说地、一起攻克难关，这些经历让我的研究生生活变得丰富多彩。还要感谢从本科相识至今的舍友浩泽、潮海和方烁。我们在求职最艰难的时刻相互扶持，在苦闷时彼此鼓励，这份友谊是我求学生涯中珍贵的财富。

上海的经历仿佛是人生中一段小小的岔路，却也遇见了许多良师益友。感谢建锋，我们仨共患难的日子应是很难忘记了。感谢方正，你对我的帮助和倾囊相授的职场智慧，让我受益良多。感谢程冉，你的激情与鼓励让我明白，永远要心存梦想。感谢於辉、鸿飞和婉婷，你们让我学会如何成为一个务实而善良的人。还要感谢朱骞、泽辰、奕宇给予我的诸多帮助，以及查琰、赵山、柯宇、刘叙、豪杰、若澎、文群、所航、梓阳、艺严、泽宇，与你们相遇是一种难得的幸运。感谢一凡、黄维、林一、居垚，是你们让我在异乡减轻了许多孤单。

感谢从高中相识至今的好友广能、起源和楚翘。从高二到现在，转眼快是十年光景。从未想过，十年之后会一起在同济的球场上挥汗如雨，一起看很多演出，一起到肇庆乡村走访调研。时光飞逝，高中时代仿佛仍在昨天，我们的友谊始终如初。

最后，将最深的感谢献给我的父母。感谢你们一路走来的鼓励与支持，你们总是竭尽所能，让我能够心无旁骛地前行。你们的爱是我最坚实的后盾。

行文至此，才发现这路上承蒙太多人的帮助与善意，寥寥数语难以道尽心中感激。今天是个浪漫的日子，我希望对所有照亮过我的人说，我爱你们。

\begin{minipage}[t]{0.945\textwidth}%
	\begin{flushright}
		2026年5月20日\\
		于华南理工大学
		\par\end{flushright}
\end{minipage}

 %致谢
\end{document}